\documentclass[12pt,a4paper]{article}
\usepackage[utf8]{inputenc}
\usepackage[T1]{fontenc}
\usepackage{amsmath,amssymb,amsthm}
\usepackage{geometry}
\usepackage{hyperref}
\usepackage{booktabs}
\usepackage{enumitem}

\geometry{margin=2.5cm}
\hypersetup{
  colorlinks=true,
  linkcolor=blue,
  urlcolor=blue,
  citecolor=blue,
  pdftitle={A Rigid Discrete Spectral Precursor from the Fibonacci Bootstrap},
  pdfauthor={Razvan-Constantin Anghelina}
}

\newtheorem{theorem}{Theorem}
\newtheorem{proposition}{Proposition}
\newtheorem{corollary}{Corollary}
\theoremstyle{definition}
\newtheorem{definition}{Definition}
\newtheorem{remark}{Remark}

\title{\textbf{A Rigid Discrete Spectral Precursor from the Fibonacci Bootstrap}}
\author{Razvan-Constantin Anghelina}
\date{April 2026}

\begin{document}
\maketitle

\begin{abstract}
We isolate an exact discrete core of the previously broader one-integer program.
Starting from the irreducible Fibonacci fusion axiom
\[
X \otimes X = \mathbf{1} \oplus X,
\]
we obtain a unique bootstrap integer
\[
a_1 = 5,
\]
and from it a rigid discrete realization centered on the regular-convex
\(H_4\) class and the 600-cell. On this finite structure we collect exact
results: the 9-point scalar spectrum in \(\mathbb{Z}[\phi]\), the McKay
\(\widetilde E_8\) shadow of the binary icosahedral group \(2I\), a discrete
Hopf layer with six compatible fibrations, a uniquely selected anisotropic wave
operator \(\Box\) on vertices, the exact kernel sector of \(\Box\), and the
exact edge-kernel decomposition
\[
\ker(\Box_1)=\rho_0 \oplus 2\rho_5.
\]

We also document two internal no-go results. First, the quotient-compatible
\(A_5\) bridge from the 12-dimensional fiber permutation module to the
nontrivial 12-dimensional edge sector is obstructed. Second, the intrinsic
\(2I\)-equivariant endomorphism algebra of the nontrivial edge sector yields a
compact Lie algebra of type \(\mathfrak{sp}(2)\), not
\(\mathfrak{u}(1)\oplus\mathfrak{su}(2)\oplus\mathfrak{su}(3)\). These two
results rule out the two gauge-completion routes tested here.

What remains is a mathematically rigid discrete spectral precursor, together
with a residual flavor scaffold consisting of three ordered family slots, a
\(\mathbb{Z}^2\) exponent lattice, and an \(A_5\)/golden-ratio background.
What does not remain is a derivation of the Standard Model gauge group, gauge
couplings, CKM/PMNS observables, or continuum gravity. The value of the present
paper is therefore twofold: an exact finite construction, and explicit no-go
results delimiting the failure of its most natural physical extensions.
\end{abstract}

\tableofcontents
\newpage

\section{Introduction}

The purpose of this paper is deliberately narrower than that of the earlier
large framework. We record only what survives a strict exact/conditional/no-go
audit. The result is not a derivation of the Standard Model. It is an exact
finite precursor together with explicit obstructions to two natural gauge
extensions.

This narrower outcome is still nontrivial. The construction produces a rigid
chain
\[
\text{Fibonacci seed} \to a_1=5 \to \mathbb{Q}(\sqrt{5}) \to H_4
\to \text{600-cell} \to \text{exact spectral data},
\]
and then yields sharp statements on the associated wave operators and kernel
sectors. Just as important, it shows where the natural physical overreads fail.

\section{Logical Status}

We separate four levels of statement.

\begin{enumerate}[label=\textup{(L\arabic*)}]
\item \textbf{Irreducible axiom.} The seed
  \[
  X \otimes X = \mathbf{1} \oplus X
  \]
  is taken as the minimal nontrivial self-reference rule.
\item \textbf{Exact discrete result.} Finite theorem or finite computation on
  the 600-cell and its associated operators.
\item \textbf{Conditional statement.} A result that holds only after importing
  an additional physical dictionary not derived here.
\item \textbf{No-go result.} A theorem or finite obstruction showing that a
  natural extension does not work.
\end{enumerate}

The present paper contains results of types (L1), (L2), and (L4), plus one
small residual flavor scaffold of type (L3).

\section{From the Fibonacci Seed to the Bootstrap Integer}

\subsection{Irreducible axiom}

The starting point is the fusion rule
\[
X \otimes X = \mathbf{1} \oplus X.
\]
In the earlier step-by-step audit, weaker principles were tested and ruled out.
The conclusion was that this rule should be treated as an irreducible axiom,
not as a theorem derived from an even weaker self-reference principle.

\subsection{Bootstrap theorem}

Let
\[
d_1(n)=\frac{\sin(3\pi/n)}{\sin(\pi/n)}.
\]
The bootstrap condition requires that the Fibonacci quantum dimension
\(\phi=(1+\sqrt{5})/2\) coincide with the level-dependent value \(d_1(n)\).

\begin{theorem}
The unique positive integer \(n\) satisfying
\[
d_1(n)=\phi
\]
is \(n=5\).
\end{theorem}

\begin{proof}[Proof sketch]
This is the exact closure of step \(S02\): a finite arithmetic verification
combined with the exact trigonometric identity at \(n=5\).
\end{proof}

We therefore write
\[
a_1=5.
\]

\section{The \(H_4\) Realization Class and the 600-Cell}

The bootstrap integer \(a_1=5\) selects the field \(\mathbb{Q}(\sqrt{5})\).
Inside the regular convex 4-polytopes, this forces attention to the \(H_4\)
pair \(120\)-cell / \(600\)-cell. The 600-cell is then singled out within that
pair by intrinsic discrete criteria: \(120\) vertices matching \(|2I|=120\)
and local vertex degree \(12\).

\begin{proposition}
Within the regular-convex \(H_4\) realization class, the 600-cell is uniquely
selected by the joint conditions
\[
|V|=120=|2I|,
\qquad
\deg(v)=12.
\]
\end{proposition}

This is a scoped selection theorem. It is not the stronger false claim that
\(a_1=5\) alone forces the 600-cell without specifying the realization class.

\section{Exact Spectral Structure}

\subsection{Scalar spectrum}

Let \(A\) be the adjacency matrix of the 600-cell graph and let
\[
\Delta_0 = 12I - A.
\]
The exact scalar spectrum contains 9 distinct eigenvalues in \(\mathbb{Z}[\phi]\),
with multiplicities
\[
1,4,9,16,25,36,9,16,4.
\]

\begin{remark}
In the audited exact-core, this is classified as a \emph{computational fact},
not as a claimed physical mass spectrum. The perfect-square multiplicities are
explained representation-theoretically and are not used here as a standalone
physical argument.
\end{remark}

\subsection{McKay shadow}

\begin{theorem}
For the binary icosahedral group \(2I\), the McKay graph of the defining
2-dimensional representation is the affine Dynkin diagram \(\widetilde E_8\).
\end{theorem}

This is a standard theorem once \(2I\) is fixed. In the present paper, the
point is not to claim that nature ``chooses \(E_8\)'', but rather that the
\(\widetilde E_8\) shadow is forced after the \(2I\) realization is chosen.

\section{Discrete Hopf Layer and Wave Operator}

\subsection{Discrete Hopf fibrations}

The relevant notion is not an arbitrary Hopf fibration of \(S^3\), but a
discrete 600-cell fibration: a partition of the 120 vertices into 12 decagons,
obtained from left cosets of a subgroup \(C_{10} \leq 2I\).

\begin{proposition}
There are exactly six discrete Hopf fibrations compatible with the \(2I\)
structure in the sense above.
\end{proposition}

\subsection{Wave operator}

For a chosen discrete Hopf fibration \(F\), let \(A_{\mathrm{fiber}}(F)\) be
the adjacency along fiber edges and write
\[
\Box_F(c)=c\,A_{\mathrm{fiber}}(F)-A.
\]

\begin{proposition}
For every one of the six compatible fibrations:
\begin{enumerate}[label=\textup{(\alph*)}]
\item the operators \(A\) and \(A_{\mathrm{fiber}}(F)\) commute;
\item the unique nontrivial integer coefficient with nonzero kernel is
  \(c=6\);
\item
  \[
  \ker(\Box_F(6))
  =
  E_A(12)\oplus E_A(6\phi)\oplus E_A(6\phi');
  \]
\item the kernel has dimension \(9\).
\end{enumerate}
\end{proposition}

The commutativity statement is structural. The uniqueness of \(c=6\) was
retained in the exact-core audit as a uniform computational fact over all six
fibrations.

\section{Edge Kernel Structure}

On the edge space, the exact-core analysis yields:
\[
\ker(\Box_1)=\rho_0 \oplus 2\rho_5.
\]
This is uniform over the six compatible discrete fibrations in the weak sense
needed by the audited chain.

Separately, the fiber permutation action factors through
\[
A_5 \cong 2I/\{\pm 1\}
\]
and the 12-dimensional permutation module on the fibers decomposes as
\[
1 \oplus 3 \oplus 3' \oplus 5.
\]

These are two distinct exact statements. What failed was the attempted bridge
between them.

\section{Two Gauge No-Go Results}

\subsection{No-go 1: the quotient-compatible \(A_5\) bridge}

The first attempted gauge route tried to use the 12-dimensional fiber
permutation module to access the nontrivial 12-dimensional edge sector.

\begin{theorem}[Gauge bridge no-go]
Any quotient-compatible \(A_5\)-equivariant map from the fiber permutation
module to \(\ker(\Box_1)\) lands only in the trivial \(\rho_0\) component.
It cannot reach the nontrivial 12-dimensional sector \(2\rho_5\).
\end{theorem}

\begin{corollary}
The fiber permutation module
\[
1 \oplus 3 \oplus 3' \oplus 5
\]
cannot serve, through this quotient-compatible route, as a derived adjoint
candidate for
\(\mathfrak{u}(1)\oplus\mathfrak{su}(2)\oplus\mathfrak{su}(3)\).
\end{corollary}

\subsection{No-go 2: the intrinsic edge-endomorphism route}

The second attempted gauge route uses only the exact nontrivial edge sector.
Let
\[
G = \ker(\Box_1)\cap \rho_0^\perp \cong 2\rho_5.
\]
One may ask whether the compact Lie algebra of \(2I\)-equivariant skew
endomorphisms of \(G\) produces the desired gauge algebra.

\begin{theorem}[Edge endomorphism no-go]
The irrep \(\rho_5\) has Frobenius--Schur indicator \(-1\), so
\[
\mathrm{End}_{2I}(\rho_5)\cong \mathbb{H},
\qquad
\mathrm{End}_{2I}(G)\cong M_2(\mathbb{H}).
\]
The associated canonical compact Lie algebra is therefore
\[
\mathfrak{sp}(2),
\]
not
\[
\mathfrak{u}(1)\oplus\mathfrak{su}(2)\oplus\mathfrak{su}(3).
\]
\end{theorem}

\subsection{Interpretation}

These two no-go theorems do not invalidate the discrete precursor. They
invalidate two specific gauge-completion routes. This distinction is the main
reason for writing the present paper separately from the earlier broader
manuscript.

\section{Residual Flavor Scaffold}

Once the gauge routes are removed, a small flavor scaffold still survives.

\subsection{Three ordered slots}

The exact-core program retains three stable chiral unit sectors. In the
flavor-first audit, these were reduced to three \emph{ordered family slots},
not physical generations in the full Standard-Model sense.

\subsection{Charge lattice}

The exact arithmetic structure is:
\[
L=\mathbb{Z}^2,
\qquad
w(a,b)=5a+6b.
\]
This gives a surjective integer grading with kernel
\[
\ker w = \mathbb{Z}(6,-5).
\]

\subsection{What remains conditional}

The following can be said, but only conditionally:

\begin{enumerate}
\item the grading \(w\) may be read as suppression-charge data;
\item an \(A_5\)/golden-ratio residual-symmetry language remains available;
\item old CKM/PMNS patterns may be reinterpreted only as candidate texture data.
\end{enumerate}

What does \emph{not} follow:

\begin{enumerate}
\item no physical mass matrix is derived;
\item no CKM matrix is derived;
\item no PMNS matrix is derived;
\item no CP observable is derived;
\item the strong-CP claim does not survive as a derived statement.
\end{enumerate}

\section{What This Paper Does Not Claim}

This paper does \emph{not} claim:

\begin{enumerate}
\item a derivation of the Standard Model gauge group;
\item a derivation of the gauge couplings \(\alpha\), \(\alpha_s\), or
  \(\sin^2\theta_W\) as physical observables;
\item a derivation of the CKM or PMNS matrices;
\item a derivation of CP observables;
\item a derivation of continuum general relativity.
\end{enumerate}

The status of these issues is now sharper than before: several natural routes
have been explicitly tested, and some of them have been ruled out.

\section{Open Problems}

The most serious open problems are now easy to state.

\begin{enumerate}
\item Find a genuinely new gauge route, if one exists at all.
\item Understand whether the residual flavor scaffold can be connected to a
  standard flavor framework without uncontrolled fitting.
\item Clarify whether the exact discrete spectral data admit a continuum shadow
  in a setting such as noncommutative geometry, discrete exterior calculus, or
  another external formalism.
\end{enumerate}

\section{Conclusion}

The final outcome is more modest than a theory of everything, but more solid.
From the Fibonacci bootstrap we obtain a rigid discrete spectral precursor:
\[
X\otimes X=\mathbf{1}\oplus X
\;\to\;
a_1=5
\;\to\;
H_4
\;\to\;
\text{600-cell}
\;\to\;
\text{exact spectra, kernels, and McKay structure}.
\]
On top of this exact core, two natural gauge-completion routes were tested and
ruled out. A small flavor scaffold survives, but not an observable-level flavor
theory.

This is therefore a paper about exact structure and explicit limits. In the
present state of the project, that is the correct claim.

\end{document}
